\bta{2019}

\begin{enumerate}
	%\renewcommand{\labelenumi}{\arabic{enumi}.}
	% A(\Alph) a(\alph) I(\Roman) i(\roman) 1(\arabic)
	%设定全局标号series=example	%引用全局变量resume=example
	%[topsep=-0.3em,parsep=-0.3em,itemsep=-0.3em,partopsep=-0.3em]
	%可使用leftmargin调整列表环境左边的空白长度 [leftmargin=0em]
	\item
选择题 (本题满分 50 分，每小题 5 分)
\begin{enumerate}
	%\renewcommand{\labelenumi}{\arabic{enumi}.}
	% A(\Alph) a(\alph) I(\Roman) i(\roman) 1(\arabic)
	%设定全局标号series=example	%引用全局变量resume=example
	%[topsep=-0.3em,parsep=-0.3em,itemsep=-0.3em,partopsep=-0.3em]
	%可使用leftmargin调整列表环境左边的空白长度 [leftmargin=0em]
	\item
求极限 \xzanswer{A}
\[
\lim \limits_{n \rightarrow \infty}\left[\left(1+\dfrac{\mathstrut 1}{\mathstrut 2!}+\dfrac{1}{3 !}+\dots+\dfrac{1}{n!}\right)+\left(\dfrac{1}{1\cdot3}+\dfrac{1}{3 \cdot 5}+\dots+\dfrac{1}{(2 n-1)(2 n+1)}\right)\right]
\]


\fourchoices
{$ e-1/2 $}
{$ 5/2 $}
{$ e+1/2 $}
{$ 7/2 $}




\item 
已知函数$ f(x) $在$ x=0 $处连续，且$\lim\limits _{h \rightarrow 0} \frac{f\left(h^{2}\right)}{h^{2}}=1$，则下列正确的是 \xzanswer{C} 


\fourchoices
{$ f(0)=0 $ 且$ f_- ^{\prime} (0) $存在}
{$ f(0)=1 $且$ f_- ^{\prime} (0) $存在}
{$ f(0)=0 $ 且$ f_+ ^{\prime} (0) $存在}
{$ f(0)=1 $ 且$ f_+ ^{\prime} (0) $存在}


\item 
以下$ 4 $个命题中正确的是 \xzanswer{D} 


\fourchoices
{若$ f ^{\prime} (x) $在$ (0,1) $内连续，则$ f(x) $在$ (0,1) $内有界}
{若$ f(x) $在$ (0,1) $内连续，则$ f(x) $在$ (0,1) $内有界}
{若$ f(x) $在$ (0,1) $内有界，则$ f ^{\prime} (x) $在$ (0,1) $内有界}
{若$ f ^{\prime} (x) $在$ (0,1) $内有界，则$ f(x) $在$ (0,1) $内有界}



\item 
设$ f(x)=\left\{\begin{array}{cl}{\dfrac{\mathstrut x \cdot \sin \left(x^{2}-1\right)}{\mathstrut x^{2}-1}} & {x \neq \pm 1} \\ \mathstrut {1} & {x=\pm 1}\end{array}\right. $，则下列说法正确的是 \xzanswer{B}


\fourchoices
{$ f(x) $在$ x=-1 $连续，在$ x=1 $不连续}
{$ f(x) $在$ x=-1 $不连续，在$ x=1 $连续}
{$ f(x) $在$ x=-1 $和$ x=1 $都连续}
{$ f(x) $在$ x=-1 $和$ x=1 $都不连续}



\item 
若向量$ \vec{a},\vec{b} $满足$ (\vec{a}+\vec{b}) \perp(\vec{a}-\vec{b}) $、$ (5\vec{a}+\vec{b}) \perp (3 \vec{a}-4 \vec{b}) $，则$\vec{a}$与$\vec{b}$的夹角$  \alpha  $满足 \xzanswer{D} 

\fourchoices
{$\cos \alpha=\frac{14}{17}$}
{$\cos \alpha=\frac{13}{17}$}
{$\cos \alpha=\frac{12}{17}$}
{$\cos \alpha=\frac{11}{17}$}







\item 
已知$ f(x, y)=\left\{\begin{array}{cl}{\dfrac{\mathstrut x^{2}+y^{2}}{\mathstrut \left|x|+y^{2}\right.},} & {x^{2}+y^{2} \neq 0} \\ \mathstrut {0} & {x^{2}+y^{2}=0}\end{array}\right. $，则关于累次极限$ I=\lim\limits _{x \rightarrow 0} \lim\limits _{y \rightarrow 0} f(x, y) $和$ J=\lim\limits _{y \rightarrow 0 } \lim\limits_{x \rightarrow 0}f(x, y) $的下列说法正确的是 \xzanswer{C} 


\fourchoices
{$I$ 存在但 $J$ 不存在}
{$J$ 存在但 $I$ 不存在}
{$I$ 和 $J$ 都存在但 $I \neq J$}
{$I$ 和 $J$ 都存在且 $I=J$}



\item 
已知正项级数$ \sum\limits_{n=0}^{\infty} a_{n} $收敛，则下列级数
$$
\sum_{n=0}^{\infty}(-1)^{n} a_{n}, \quad \sum_{n=0}^{\infty}\left(a_{n}\right)^{2}, \quad  \sum_{n=0}^{\infty} \sin \left(a_{n}\right),  \quad \sum_{n=0}^{\infty} \ln \left(1+a_{n}\right)
$$
中收敛级数的数量是 \xzanswer{D}


\fourchoices
{$ 1 $}
{$ 2 $}
{$ 3 $}
{$ 4 $}

\item 
设函数$ f(t) $二次连续可微，令$ u=f(xy) $，则$ \phi(t)=\dfrac{\partial^{2} u}{\partial x \partial y} $的表达式是 \xzanswer{C}


\fourchoices
{$t f^{\prime \prime}(t)-f^{\prime}(t)$}
{$f^{\prime \prime}(t)-t f^{\prime}(t)$}
{$ t f^{\prime \prime}(t)+f^{\prime}(t)$}
{$ f^{\prime \prime}(t)+t f^{\prime}(t)$}


\item 
设非齐次线性微分方程$ y ^{\prime} +P(x)y=Q(x) $有两个不同的解$y_{1}(x), y_{2}(x)$，$ C $为任意常数，则该方程的通解是 \xzanswer{B} 


\fourchoices
{$C\left[y_{1}(x)-y_{2}(x)\right]$}
{$y_{1}(x)+C\left[y_{1}(x)-y_{2}(x)\right]$}
{$C\left[y_{1}(x)+y_{2}(x)\right]$}
{$y_{1}(x)+C\left[y_{1}(x)+y_{2}(x)\right]$}

\item 
考虑积分$ I=\iint_{D}\left(x^{2}-y^{2}\right) d x d y $，其中$ D $是椭圆：$ \{(x, y) \in R^{2} | \dfrac{x^{2}}{a^{2}}+\dfrac{y^{2}}{b^{2}}<1\} $落在第一象限中的部分。则下列说法正确的是 \xzanswer{A}


\fourchoices
{$ a>b $时$ I>0 $}
{$ a>b $时$ I<0 $}
{$ a>b $时$ I=0 $}
{$ a<b $时$ I=0 $}




\end{enumerate}


\banswer{
\begin{enumerate}
	%\renewcommand{\labelenumi}{\arabic{enumi}.}
	% A(\Alph) a(\alph) I(\Roman) i(\roman) 1(\arabic)
	\item
	\textbf{A}。两个括号分别求极限。由 $e=\sum\limits_{k=0}^{\infty}\frac{1}{k!}=1+\left(1+\frac{1}{2!}+\cdots\right)$ 得
	\[
	1+\frac{1}{2!}+\cdots+\frac{1}{n!}\longrightarrow e-1 ;
	\]
	又 $\frac{1}{(2k-1)(2k+1)}=\frac12\left(\frac{1}{2k-1}-\frac{1}{2k+1}\right)$，裂项求和得
	$\frac12\left(1-\frac{1}{2n+1}\right)\to\frac12$。故原极限 $=e-1+\frac12=e-\frac12$。

	\item
	\textbf{C}。$h\to0$ 时 $h^{2}\to0^{+}$，题设即 $\lim\limits_{u\to0^{+}}\frac{f(u)}{u}=1$。
	分母趋于 $0$ 而比值有有限极限，故 $f(u)\to0\ (u\to0^{+})$；再由 $f$ 在 $x=0$ 连续得 $f(0)=0$。
	于是 $f_{+}^{\prime}(0)=\lim\limits_{u\to0^{+}}\frac{f(u)-f(0)}{u}=1$ 存在。
	因 $h^{2}>0$，题设只给出右侧信息，$f_{-}^{\prime}(0)$ 未必存在。

	\item
	\textbf{D}。若 $\left|f^{\prime}\right|\leqslant M$ 于 $(0,1)$，则由拉格朗日中值定理
	$\left|f(x)-f\left(\frac12\right)\right|\leqslant M\left|x-\frac12\right|\leqslant M$，故 $f$ 有界。
	A、B 的反例都可取 $f(x)=\frac1x$（在 $(0,1)$ 内连续可微而无界）；
	C 的反例取 $f(x)=\sin\frac1x$（有界，但 $f^{\prime}=-\frac{1}{x^{2}}\cos\frac1x$ 无界）。

	\item
	\textbf{B}。$x\to\pm1$ 时 $x^{2}-1\to0$，故 $\frac{\sin\left(x^{2}-1\right)}{x^{2}-1}\to1$，从而
	$\lim\limits_{x\to1}f(x)=1\cdot1=1=f(1)$，在 $x=1$ 连续；而
	$\lim\limits_{x\to-1}f(x)=(-1)\cdot1=-1\neq1=f(-1)$，在 $x=-1$ 不连续。

	\item
	\textbf{D}。由 $(\vec{a}+\vec{b})\perp(\vec{a}-\vec{b})$ 得
	$|\vec{a}|^{2}-|\vec{b}|^{2}=0$，记 $m^{2}=|\vec{a}|^{2}=|\vec{b}|^{2}$。
	再由 $(5\vec{a}+\vec{b})\perp(3\vec{a}-4\vec{b})$ 得
	\[
	15|\vec{a}|^{2}-17\,\vec{a}\cdot\vec{b}-4|\vec{b}|^{2}=0
	\;\Longrightarrow\;
	11m^{2}=17\,\vec{a}\cdot\vec{b},
	\]
	故 $\cos\alpha=\frac{\vec{a}\cdot\vec{b}}{|\vec{a}||\vec{b}|}=\frac{11}{17}$。

	\item
	\textbf{C}。固定 $x\neq0$ 令 $y\to0$ 得 $\frac{x^{2}}{|x|}=|x|$，再令 $x\to0$ 得 $I=0$；
	固定 $y\neq0$ 令 $x\to0$ 得 $\frac{y^{2}}{y^{2}}=1$，再令 $y\to0$ 得 $J=1$。
	两个累次极限都存在但 $I\neq J$（于是重极限必不存在）。

	\item
	\textbf{D}。由 $a_n>0$ 且 $\sum a_n$ 收敛知 $a_n\to0$，故存在 $N$ 使 $n>N$ 时 $0<a_n<1$。
	此时 $\left|(-1)^{n}a_n\right|=a_n$、$a_n^{2}<a_n$、$0<\sin a_n\leqslant a_n$、
	$0<\ln\left(1+a_n\right)\leqslant a_n$，由比较判别法四个级数全部收敛，故个数为 $4$。

	\item
	\textbf{C}。$\frac{\partial u}{\partial x}=yf^{\prime}(xy)$，再对 $y$ 求导得
	\[
	\frac{\partial^{2}u}{\partial x\partial y}=f^{\prime}(xy)+xy\,f^{\prime\prime}(xy) ,
	\]
	令 $t=xy$ 即 $\phi(t)=tf^{\prime\prime}(t)+f^{\prime}(t)$。

	\item
	\textbf{B}。$y_1-y_2$ 满足对应齐次方程 $y^{\prime}+P(x)y=0$ 且不恒为零，
	而 $y_1$ 是非齐次方程的一个特解，由“非齐次通解 $=$ 一个特解 $+$ 齐次通解”得
	$y=y_1+C\left(y_1-y_2\right)$。

	\item
	\textbf{A}。作广义极坐标 $x=ar\cos\theta$、$y=br\sin\theta$，Jacobi 行列式为 $abr$：
	\begin{align*}
	I&=ab\int_0^{\pi/2}\!\!\int_0^1\left(a^{2}\cos^{2}\theta-b^{2}\sin^{2}\theta\right)r^{3}dr\,d\theta\\
	&=\frac{ab}{4}\cdot\frac{\pi}{4}\left(a^{2}-b^{2}\right)=\frac{\pi ab\left(a^{2}-b^{2}\right)}{16} .
	\end{align*}
	故 $a>b$ 时 $I>0$（$a<b$ 时 $I<0$，$a=b$ 时才为 $0$）。
\end{enumerate}

	\par\vspace{1mm}\noindent
	\textbf{注：}{\kaishu
	第 (7) 题题面原作“已知数项级数”，漏掉了“正项”二字，已据题意补上；原答案 A 有误。
	若只假定 $\sum a_n$ 收敛（可以条件收敛），则四个级数无一必收敛，四个选项都不正确：
	取 $a_n=\frac{(-1)^{n}}{n}$ 则 $\sum(-1)^{n}a_n=\sum\frac1n$ 发散；
	取 $a_n=\frac{(-1)^{n}}{\sqrt n}$ 则 $\sum a_n^{2}=\sum\frac1n$ 及
	$\sum\ln\left(1+a_n\right)$ 都发散；
	把各项按 $\frac{1}{k^{1/3}},\ \frac{1}{k^{1/3}},\ -\frac{2}{k^{1/3}}$ 三个一组排成 $\left\{a_n\right\}$，
	则 $\sum a_n$ 收敛而 $\sum\sin a_n$ 发散（因 $\sin a\approx a-\frac{a^{3}}{6}$，
	而 $\sum a_n^{3}$ 的通项相当于 $-\frac6k$）。以上均经数值计算验证。
	$2022$ 年试卷中的同一道题写明“正项级数”，答案正是 D。
	}\par
}


\item
(本题满分 10 分)
求常微分方程：$yy''+(y')^2=0$满足初始条件：$y(0)=1,y'(0)= \frac{ 1 }{ 2 } $的特解。


\banswer{
	注意 $yy^{\prime\prime}+\left(y^{\prime}\right)^{2}=\left(yy^{\prime}\right)^{\prime}$，
	故方程即 $\left(yy^{\prime}\right)^{\prime}=0$，于是 $yy^{\prime}\equiv C_1$。
	由 $y(0)=1$、$y^{\prime}(0)=\frac12$ 得 $C_1=\frac12$。

	又 $yy^{\prime}=\left(\frac{y^{2}}{2}\right)^{\prime}$，故 $\frac{y^{2}}{2}=\frac{x}{2}+C_2$，
	即 $y^{2}=x+C_3$；代入 $y(0)=1$ 得 $C_3=1$。因 $y(0)=1>0$，取正根，所求特解为
	\[
	y^{2}=x+1,\qquad\text{即}\quad y=\sqrt{x+1}\quad(x>-1) .
	\]
	（验证：$y^{\prime}=\frac{1}{2\sqrt{x+1}}$，$y^{\prime}(0)=\frac12$；
	$yy^{\prime}\equiv\frac12$ 为常数，故 $yy^{\prime\prime}+\left(y^{\prime}\right)^{2}=0$。）
}


\item 
(本题满分 10 分)
求直线$ L $：$ \frac{x-1}{2}=\frac{y+1}{-1}=\frac{z}{1} $绕$ z $轴旋转一周所生成的旋转曲面的方程。


\banswer{
	先把 $L$ 化为以 $z$ 为参数的形式。令
	$\frac{x-1}{2}=\frac{y+1}{-1}=\frac{z}{1}=t$，得
	\[
	x=1+2t,\qquad y=-1-t,\qquad z=t ,
	\]
	即 $t=z$ 时 $x=1+2z$、$y=-1-z$。

	绕 $z$ 轴旋转时，$L$ 上高度为 $z$ 的点扫出以 $z$ 轴为心、半径为该点到 $z$ 轴距离的圆，
	故旋转面上的点 $(x,y,z)$ 满足
	\[
	x^{2}+y^{2}=(1+2z)^{2}+(-1-z)^{2}=(1+2z)^{2}+(1+z)^{2} ,
	\]
	展开即
	\[
	x^{2}+y^{2}=5z^{2}+6z+2 .
	\]
	配方可写成
	\[
	\frac{x^{2}+y^{2}}{1/5}-\frac{\left(z+\frac35\right)^{2}}{1/25}=1 ,
	\]
	是一个单叶旋转双曲面（因为 $L$ 与 $z$ 轴异面）。
}


\item 
(本题满分 10 分)
设$ a_0=3,a_1=5 $，且对任何自然数$ n>1 $有$ n a_{n}=\frac{2}{3} a_{n-1}-\left(n-1\right)a_{n-1} $。
证明：当$ |x|<1 $时，幂级数$ \sum\limits_{n=0}^{\infty} a_{n} x^{n} $
绝对收敛，并求其和函数$ S(x) $。


\banswer{
	\textbf{绝对收敛.} 递推式即 $a_n=\frac{\frac53-n}{n}a_{n-1}\ (n\geqslant2)$，
	由 $a_1=5\neq0$ 知一切 $a_n\neq0$，且
	\[
	\left|\frac{a_nx^{n}}{a_{n-1}x^{n-1}}\right|=\frac{\left|n-\frac53\right|}{n}\,|x|
	\longrightarrow|x| .
	\]
	由比值判别法，$|x|<1$ 时 $\sum a_nx^{n}$ 绝对收敛（$|x|>1$ 时发散，故收敛半径为 $1$）。

	\textbf{求和函数.} 设 $S(x)=\sum\limits_{n=0}^{\infty}a_nx^{n}$，$|x|<1$。
	把 $na_n=\left(\frac23-n+1\right)a_{n-1}$ 两端乘 $x^{n-1}$ 并对 $n\geqslant2$ 求和：
	左端为 $\sum\limits_{n\geqslant2}na_nx^{n-1}=S^{\prime}(x)-a_1$；
	右端令 $m=n-1$ 得
	\[
	\sum_{m\geqslant1}\left(\frac23-m\right)a_mx^{m}
	=\frac23\left[S(x)-a_0\right]-xS^{\prime}(x) .
	\]
	于是
	\[
	(1+x)S^{\prime}=\frac23S-\frac23a_0+a_1=\frac23S+3,\qquad S(0)=a_0=3 .
	\]
	分离变量
	\[
	\frac{dS}{\frac23S+3}=\frac{dx}{1+x}
	\;\Longrightarrow\;
	\frac23S+3=K(1+x)^{\frac23},
	\]
	代入 $x=0$ 得 $K=5$。故
	\[
	S(x)=\frac{15}{2}(1+x)^{\frac{2}{3}}-\frac{9}{2},\qquad |x|<1 .
	\]
	（校验：$(1+x)^{2/3}=1+\frac23x-\frac19x^{2}+\cdots$，代入得
	$S=3+5x-\frac56x^{2}+\cdots$，与 $a_0=3$、$a_1=5$、$a_2=-\frac56$ 相符；
	取 $x=0.3,\,-0.5,\,0.8$ 数值求和亦与该闭式一致。）
}


\item 
(本题满分 10 分)
已知函数$ y=y(x) $由方程：$ e^{y}+6 x y+x^{2}-1=0 $确定，试求$ y^{\prime\prime} (0) $。

\banswer{
	\textbf{定 $y(0)$.} 在 $e^{y}+6xy+x^{2}-1=0$ 中令 $x=0$ 得 $e^{y}=1$，故 $y(0)=0$。

	\textbf{定 $y^{\prime}(0)$.} 两端对 $x$ 求导：
	\[
	e^{y}y^{\prime}+6y+6xy^{\prime}+2x=0 .
	\]
	代入 $x=0$、$y=0$ 得 $y^{\prime}(0)=0$。

	\textbf{定 $y^{\prime\prime}(0)$.} 再对 $x$ 求导一次：
	\[
	e^{y}\left(y^{\prime}\right)^{2}+e^{y}y^{\prime\prime}+12y^{\prime}+6xy^{\prime\prime}+2=0 .
	\]
	代入 $x=0$、$y=0$、$y^{\prime}(0)=0$ 得 $y^{\prime\prime}(0)+2=0$，即
	\[
	y^{\prime\prime}(0)=-2 .
	\]
}

\newpage
\item 
(本题满分 10 分)
对实数$ R>0 $定义积分：
\[ 
I_R=\iiint\limits_{1/R \leq x^2+y^2+z^2 \leq R}\dfrac{e^{-(x^2+y^2+z^2)}}{x^2+y^2+z^2}dxdydz
 \]
证明极限$ I=\lim \limits _{R \rightarrow+\infty} I_{R} $存在并计算其值。

\banswer{
	用球坐标 $x=\rho\sin\varphi\cos\theta$ 等。积分区域 $\frac1R\leqslant\rho^{2}\leqslant R$ 即
	$\frac{1}{\sqrt R}\leqslant\rho\leqslant\sqrt R$（须 $R>1$），
	被积函数只依赖 $\rho$，且体积元中的 $\rho^{2}$ 恰与分母约去：
	\begin{align*}
	I_R&=\int_0^{2\pi}\!\!d\theta\int_0^{\pi}\!\sin\varphi\,d\varphi
	\int_{1/\sqrt R}^{\sqrt R}\frac{e^{-\rho^{2}}}{\rho^{2}}\,\rho^{2}d\rho\\
	&=4\pi\int_{1/\sqrt R}^{\sqrt R}e^{-\rho^{2}}d\rho .
	\end{align*}

	\textbf{极限存在.} $R$ 增大时区间 $\left[\frac{1}{\sqrt R},\sqrt R\right]$ 单调扩张，
	被积函数 $e^{-\rho^{2}}>0$，故 $I_R$ 关于 $R$ 单调增加；又
	\[
	I_R\leqslant4\pi\int_0^{+\infty}e^{-\rho^{2}}d\rho<+\infty
	\]
	有上界。单调有界必有极限，故 $I=\lim\limits_{R\to+\infty}I_R$ 存在。

	\textbf{求值.} 由 Poisson 积分 $\int_0^{+\infty}e^{-\rho^{2}}d\rho=\frac{\sqrt\pi}{2}$，
	\[
	I=4\pi\cdot\frac{\sqrt{\pi}}{2}=2\pi^{3/2}\approx11.137 .
	\]
}


\item 
(本题满分 10 分)
计算积分
\[ 
I=\oint\limits_{L} \frac{\left(x^{2}-y^{2}-x\right) d y+(1-2 x) y d x}{\left(x^{2}+y^{2}\right)\left[(x-1)^{2}+y^{2}\right]}
 \]
其中$ L=\left\{(x, y) \in R^{2} \big| x^{2}+y^{2}=4\right\} $沿逆时针方向。

\banswer{
	记
	\[
	u=x^{2}+y^{2},\qquad v=(x-1)^{2}+y^{2}=u-2x+1 ,
	\]
	则 $v-u=1-2x$。关键是把被积表达式拆成两个“标准”角微分之和。

	\textbf{拆项.} 分子中 $(1-2x)y=(v-u)y$，故
	\[
	P:=\frac{(1-2x)y}{uv}=\frac{y}{u}-\frac{y}{v} .
	\]
	又由 $xv-(x-1)u=x(u-2x+1)-(xu-u)=-\left(x^{2}-y^{2}-x\right)$ 得
	\[
	Q:=\frac{x^{2}-y^{2}-x}{uv}=\frac{x-1}{v}-\frac{x}{u} .
	\]
	于是
	\begin{align*}
	P\,dx+Q\,dy
	&=-\frac{-y\,dx+x\,dy}{x^{2}+y^{2}}
	+\frac{-y\,dx+(x-1)dy}{(x-1)^{2}+y^{2}} .
	\end{align*}

	\textbf{逐项积分.} 两个奇点 $(0,0)$ 与 $(1,0)$ 都在 $L:\ x^{2}+y^{2}=4$ 的内部。
	上式每一项在各自奇点之外都满足 $\frac{\partial Q}{\partial x}=\frac{\partial P}{\partial y}$，
	故可用挖点法把 $L$ 分别收缩到绕该奇点的小圆周，得到熟知的结果
	\[
	\oint_L\frac{-y\,dx+x\,dy}{x^{2}+y^{2}}=2\pi,\qquad
	\oint_L\frac{-y\,dx+(x-1)dy}{(x-1)^{2}+y^{2}}=2\pi .
	\]
	两者相减，
	\[
	I=-2\pi+2\pi=0 .
	\]
	（沿 $x=2\cos t$、$y=2\sin t$ 作数值积分，结果确为 $0$。）
}


\item 
(本题满分 10 分)
已知$ f(x)=\sin \frac{1}{x} $，试求一个序列$ x_{n} \rightarrow 0 \ (n \rightarrow \infty) $使得$ \lim\limits _{n \rightarrow \infty} f\left(x_{n}\right) $存在，且满足
$ \lim\limits _{x \rightarrow 0}\left[\lim\limits _{n \rightarrow \infty} f\left(x_{n}\right)+x\right]^{\left[ 2/f(\frac{1}{x}) \right]}=e^{2} $

\banswer{
	\textbf{化简指数.} 注意 $f\left(\frac1x\right)=\sin\frac{1}{1/x}=\sin x$，
	故所给极限就是 $\lim\limits_{x\to0}(A+x)^{\frac{2}{\sin x}}$，
	其中 $A=\lim\limits_{n\to\infty}f\left(x_n\right)$。

	\textbf{定 $A$.} 指数 $\frac{2}{\sin x}\to\infty$。
	若 $0<A<1$ 则极限为 $0$；若 $A>1$ 则极限为 $+\infty$；都不等于 $e^{2}$。
	故必须 $A=1$，此时为 $1^{\infty}$ 型：
	\[
	\ln\left[(1+x)^{\frac{2}{\sin x}}\right]=\frac{2\ln(1+x)}{\sin x}\sim\frac{2x}{x}=2 ,
	\]
	极限恰为 $e^{2}$，符合要求。

	\textbf{构造 $x_n$.} 只需 $x_n\to0$ 且 $\sin\frac{1}{x_n}\to1$。
	取 $\frac{1}{x_n}=2n\pi+\frac{\pi}{2}$，即
	\[
	x_n=\frac{1}{2 n\pi +\frac{1}{2} \pi}\longrightarrow0,\qquad
	f\left(x_n\right)=\sin\left(2n\pi+\frac{\pi}{2}\right)=1 ,
	\]
	于是 $\lim\limits_{n\to\infty}f\left(x_n\right)=1=A$。答案不唯一：
	任何满足 $x_n\to0$ 且 $\sin\frac{1}{x_n}\to1$ 的序列都可以。
}


\item 
(本题满分 10 分)
设函数$ f(x) $在$[a, b]$上$ 2 $阶可导，其中$ a<b $。且有$ f^{\prime}\left(\frac{a+b}{2}\right)=0 $。
证明：存在$ \xi	\in(a,b) $使得$\left|f^{\prime \prime}(\xi)\right| \geq \frac{4}{(b-a)^{2}}|f(b)-f(a)|$。


\banswer{
	记 $x_0=\frac{a+b}{2}$。由 Taylor 公式，将 $f(a)$ 与 $f(b)$ 在 $x_0$ 处展开：
	\[
	f(a)=f(x_0)+f'(x_0)(a-x_0)+\frac{f''(\xi_1)}{2}(a-x_0)^2,\quad \xi_1\in(a,x_0)
	\]
	\[
	f(b)=f(x_0)+f'(x_0)(b-x_0)+\frac{f''(\xi_2)}{2}(b-x_0)^2,\quad \xi_2\in(x_0,b)
	\]
	由 $f'(x_0)=0$ 及 $a-x_0=-\frac{b-a}{2}$，$b-x_0=\frac{b-a}{2}$，得
	\[
	f(b)-f(a)=\frac{(b-a)^2}{8}\bigl[f''(\xi_2)-f''(\xi_1)\bigr].
	\]
	于是
	\[
	|f(b)-f(a)|\leq\frac{(b-a)^2}{8}\bigl(|f''(\xi_1)|+|f''(\xi_2)|\bigr).
	\]
	取 $|f''(\xi)|=\max\{|f''(\xi_1)|,|f''(\xi_2)|\}$，则
	\[
	|f(b)-f(a)|\leq\frac{(b-a)^2}{8}\cdot2|f''(\xi)|=\frac{(b-a)^2}{4}|f''(\xi)|,
	\]
	即 $|f''(\xi)|\geq\frac{4}{(b-a)^2}|f(b)-f(a)|$，其中 $\xi\in(a,b)$。
}


\item 
(本题满分 10 分)
证明：
\begin{enumerate}
	%\renewcommand{\labelenumi}{\arabic{enumi}.}
	% A(\Alph) a(\alph) I(\Roman) i(\roman) 1(\arabic)
	%设定全局标号series=example	%引用全局变量resume=example
	%[topsep=-0.3em,parsep=-0.3em,itemsep=-0.3em,partopsep=-0.3em]
	%可使用leftmargin调整列表环境左边的空白长度 [leftmargin=0em]
	\item
闭区间$ [0,1] $上任何连续函数都有原函数。

\item
闭区间$ [0,1] $上任何连续可微函数都可以写成两个单调不减函数之差。
\end{enumerate}

\banswer{
\begin{enumerate}
	%\renewcommand{\labelenumi}{\arabic{enumi}.}
	% A(\Alph) a(\alph) I(\Roman) i(\roman) 1(\arabic)
	%设定全局标号series=example	%引用全局变量resume=example
	%[topsep=-0.3em,parsep=-0.3em,itemsep=-0.3em,partopsep=-0.3em]
	%可使用leftmargin调整列表环境左边的空白长度 [leftmargin=0em]
	\item
	设 $f\in C[0,1]$，令 $F(x)=\int_0^{x}f(t)\,dt$。
	由微积分基本定理（变上限积分求导），$f$ 在 $[0,1]$ 上连续蕴含 $F$ 在 $[0,1]$ 上可导且
	\[
	F^{\prime}(x)=f(x),\qquad x\in[0,1]
	\]
	（端点理解为单侧导数）。故 $F$ 就是 $f$ 的一个原函数，即任何连续函数都有原函数。
	\item 
	设 $f\in C^{1}[0,1]$。把 $f^{\prime}$ 分解为正部与负部：
	\[
	p(t)=\frac{\left|f^{\prime}(t)\right|+f^{\prime}(t)}{2},\qquad
	q(t)=\frac{\left|f^{\prime}(t)\right|-f^{\prime}(t)}{2},
	\]
	则 $p,q$ 连续、非负，且 $p-q=f^{\prime}$。令
	\[
	g(x)=f(0)+\int_0^{x}p(t)\,dt,\qquad h(x)=\int_0^{x}q(t)\,dt .
	\]
	因被积函数非负，$g$、$h$ 都在 $[0,1]$ 上单调不减；又
	\[
	g(x)-h(x)=f(0)+\int_0^{x}f^{\prime}(t)\,dt=f(x) .
	\]
	故 $f=g-h$ 正是两个单调不减函数之差。

	\par\vspace{1mm}\noindent
	\textbf{注：}{\kaishu
	原答案对第 (2) 问标为“存疑”，其实结论成立，上面的正负部分解即是证明。
	更一般地，$[0,1]$ 上任何有界变差函数都可写成两个单调不减函数之差（Jordan 分解），
	连续可微函数只是其中的特例。
	}\par
\end{enumerate}

	
}


\item 
(本题满分 10 分)
求函数$ f(x, y, z)=\ln x+2 \ln y+3 \ln z $在球面$ x^{2}+y^{2}+z^{2}=6 r^{2} $上的最大值，并证明对任意正实数$ a,b,c $，不等式$ a b^{2} c^{3} \leq 108\left(\frac{a+b+c}{6}\right)^6 $恒成立。

\banswer{
	\textbf{第一步：求条件最大值.} 在 $x,y,z>0$、$x^{2}+y^{2}+z^{2}=6r^{2}$ 上用 Lagrange 乘数法。令
	\[
	L=\ln x+2\ln y+3\ln z-\lambda\left(x^{2}+y^{2}+z^{2}-6r^{2}\right),
	\]
	由 $\frac1x=2\lambda x$、$\frac2y=2\lambda y$、$\frac3z=2\lambda z$ 得
	$2\lambda x^{2}=1$、$2\lambda y^{2}=2$、$2\lambda z^{2}=3$，即
	\[
	x^{2}:y^{2}:z^{2}=1:2:3 .
	\]
	与约束联立得 $x^{2}=r^{2}$、$y^{2}=2r^{2}$、$z^{2}=3r^{2}$，唯一驻点为
	$\left(r,\sqrt2\,r,\sqrt3\,r\right)$。当 $x,y,z$ 中有一个趋于 $0$ 时 $f\to-\infty$，
	故该驻点给出最大值
	\begin{align*}
	f_{\max}&=\ln r+2\ln\left(\sqrt2\,r\right)+3\ln\left(\sqrt3\,r\right)\\
	&=\ln\left(r\cdot2r^{2}\cdot3\sqrt3\,r^{3}\right)=\ln\left(6\sqrt3\,r^{6}\right).
	\end{align*}
	等价地：$x,y,z>0$ 且 $x^{2}+y^{2}+z^{2}=6r^{2}$ 时
	\[
	x y^{2} z^{3}\leqslant6\sqrt{3}\,r^{6} .
	\]

	\textbf{第二步：证明不等式.} 对任意正实数 $a,b,c$，取
	\[
	x=\sqrt a,\quad y=\sqrt b,\quad z=\sqrt c,\qquad r^{2}=\frac{a+b+c}{6} ,
	\]
	则 $x^{2}+y^{2}+z^{2}=a+b+c=6r^{2}$，即点 $\left(\sqrt a,\sqrt b,\sqrt c\right)$ 恰落在球面上。
	代入第一步的结论，
	\[
	\sqrt a\cdot b\cdot c^{\frac32}\leqslant6\sqrt3\,r^{6}
	=6\sqrt{3}\left(\frac{a+b+c}{6}\right)^{3} .
	\]
	两端均为正，平方即得（注意 $\left(6\sqrt3\right)^{2}=108$）
	\[
	a b^{2} c^{3}\leqslant108\left(\frac{a+b+c}{6}\right)^{6} ,
	\]
	等号当且仅当 $a:b:c=1:2:3$ 时成立。
}
	
	
	
\end{enumerate}

